\documentclass[12pt]{article}

\usepackage[portuguese]{babel}
\usepackage{float}
\usepackage{tikz}
\usepackage{listings}
\usepackage{indentfirst}
\usepackage{amsfonts, amssymb, amsmath}
\usepackage{graphicx}
\usepackage{enumitem}
\usepackage[colorlinks=true, allcolors=blue]{hyperref}
\usepackage[letterpaper,
            top=1.5cm,
            bottom=1.5cm,
            left=2.75cm,
            right=2.75cm,
            marginparwidth=1.75cm]{geometry}
\usetikzlibrary{calc,matrix}

\title{Pesquisa sobre Matrizes e Determinantes}
\author{Guilherme dos Santos Lira Lopes RA: 143630}


\begin{document}
\maketitle

\vspace{16pt}

Matrizes e determinantes são conceitos básicos da matemática, com usos variados 
e de extrema importância na computação, com amplas áreas que utilizam deles, 
como na cibersegurança com a criptografia, no design gráfico com o uso de 
operações em matrizes para rotacionar imagens além de ser a base fundamental 
para a IA. Além da computação, matrizes e determinantes tem usos na engenharia 
para analise de estruturas, circuitos elétricos e calculos de tensão e em 
estatistica para análises de dados e modelos de previsão


\section{Matrizes}
Matrizes são conjuntos de elementos dispostos tabelas, onde $i$ representa 
as linhas da tabela e $j$ representa as colunas. Uma matriz $A$ de ordem 
$m\times n$ pode ser indicada por: $A_{m\times n} = [a_{ij}]_{m\times n}$, 
com $i = \{1, 2, 3, \ldots, m\}$ e $j = \{1, 2, 3, \ldots, n\}$, com $a_{ij}$
sendo o elemento referente à linha $i$ e coluna $j$.

Uma matriz $A_{m\times n}$ pode ser representada das seguinte forma:
\[
\begin{pmatrix}
    a_{11} & a_{12} & a_{13} & \cdots & a_{1n} \\
    a_{21} & a_{22} & a_{23} & \cdots & a_{2n} \\
    a_{31} & a_{32} & a_{33} & \cdots & a_{3n} \\
    \vdots & \vdots & \vdots & \ddots & \vdots \\
    a_{m1} & a_{m2} & a_{m3} & \cdots & a_{mn}
\end{pmatrix}
\]

Uma matriz pode possuir diversas formas, podendo ser:
\begin{itemize}
    \item Matriz unidade, onde $m=n=1$
    \item Matriz linha ou coluna, onde $m$ ou $n$ são iguais a 1
    \item Matriz quadrada, onde $m=n$
\end{itemize} 

\vspace{16pt}

Dentro das matrizes quadradas, também temos alguns conceitos importantes:
\begin{itemize}
    \item Matriz unitária, onde $A_{1 \times 1}$
    \[
    \begin{pmatrix}
        a_{11}
    \end{pmatrix}
    \]
    \item Diagonal principal, formada pelos elementos $a_{ij}$, onde $i = j$
    \item Diagonal secundária, formada pelos elementos $a_{ij}$, onde 
    $i+j = m + 1$ 
    \item Matriz identidade, onde os elementos da diagonal principal são 
    iguais a 1, e todos os outros iguais a 0.
    \[
    \begin{pmatrix}
        1 & 0 & 0 \\
        0 & 1 & 0 \\
        0 & 0 & 1 \\
    \end{pmatrix}
    \]
    \item Matrizes triangulares -- Podemos ter dois tipos de matrizes 
    triangulares:
    \begin{enumerate}
        \item Matriz triangular superior -- Onde os elementos abaixo da 
        diagonal principal são iguais a 0
        \[
        A_3 = \left\{
        \begin{array}{l}
            a_{ij} = 0 \rightarrow i > j \\
            a_{ij} \in \mathbb{R} \rightarrow i \leq j
        \end{array} \right.
        \]
        \[ A_3 = 
        \begin{pmatrix}
            4 & 9 & 7 \\
            0 & 2 & 6 \\
            0 & 0 & 5 \\
        \end{pmatrix}
        \]
        \item Matriz triangular inferior -- Onde os elementos acima da 
        diagonal principal são iguais a 0
        \[
        A_3 = \left\{
        \begin{array}{l}
            a_{ij} = 0 \rightarrow i < j \\
            a_{ij} \in \mathbb{R} \rightarrow i \geq j
        \end{array} \right.
        \]
        \[ A_3 = 
        \begin{pmatrix}
            4 & 0 & 0 \\
            7 & 1 & 0 \\
            5 & 3 & 5 \\
        \end{pmatrix}
        \]
    \end{enumerate}
\end{itemize} 

\section{Determinantes}
O determinante é um número real associado exclusivamente a matrizes quadradas
que pode ter diversos significados dependendo do contexto. Determinantes 
possuem diversos calculos especificos para diversos casos de matrizes.

O determinante de uma matriz $A_m$ pode ser representado da seguinte maneira:
\[
\begin{vmatrix}
    a_{11} & a_{12} & a_{13} & \cdots & a_{1n} \\
    a_{21} & a_{22} & a_{23} & \cdots & a_{2n} \\
    a_{31} & a_{32} & a_{33} & \cdots & a_{3n} \\
    \vdots & \vdots & \vdots & \ddots & \vdots \\
    a_{m1} & a_{m2} & a_{m3} & \cdots & a_{mn}
\end{vmatrix}
\]

\begin{itemize}
    \item Determinante de matriz de ordem 1 ($A_1$):
    \[ \begin{vmatrix} a_{1} \end{vmatrix} = 1 \]
    \item Determinante de matriz de ordem 2 ($A_2$):
    \begin{itemize}
        \item O determinante de uma matriz $A_2$ é o produto dos elementos 
        da diagonal principal, menos o produto dos elementos da diagonal 
        secundária
    \end{itemize}
    \[
    \begin{vmatrix}
        a_{11} & a_{12} \\
        a_{21} & a_{22} 
    \end{vmatrix}
    = a_{11} \cdot a_{22} - (a_{12} \cdot a_{21})
    \]
    \item Determinantes de matrizes triangulares:
    \begin{itemize}
        \item Independente qual seja, superior ou inferior, o determinante 
        é sempre o produto da diagonal principal
        \[
        \begin{vmatrix}
            4 & 0 & 0 \\
            6 & 1 & 0 \\
            10 & 891 & 5 \\
        \end{vmatrix}
        = 4 \cdot 1 \cdot 5 = 20
        \] 
        \item Determinante de matriz de ordem a partir de 3
        \begin{itemize}
            \item Existem diversas formas de calcular o determinante de 
            uma matriz $3\times3$, $4\times4$, etc.
            \item Para matrizes de ordem 3 temos a Regra de Sarrus, e para 
            matrizes com ordem maior que 3, podemos usar o Teorema de Laplace 
        \end{itemize}
    \end{itemize}
\end{itemize}
    \subsection{Regra de Sarrus}
    \begin{itemize}
        \item Essa regra se aplica somente a matrizes de ordem 3, e consistem em:
        \begin{enumerate}
            \item Duplicar as duas primeiras colunas da matriz:
            \[
            \begin{vmatrix}
                a_{11} & a_{12} & a_{13} \\
                a_{21} & a_{22} & a_{23} \\
                a_{31} & a_{32} & a_{33} \\
            \end{vmatrix}
            \begin{matrix}
                a_{11} & a_{12} \\
                a_{21} & a_{22} \\
                a_{31} & a_{32}
            \end{matrix}
            \]
            \item Multiplicar os elementos das diagonais principais e diagonais 
            secundárias:
            \[
            \begin{tikzpicture}[>=stealth]
                \matrix [%
                matrix of math nodes,
                column sep=1em,
                row sep=1em
                ] (sarrus) {%
                a_{11} & a_{12} & a_{13} & a_{11} & a_{12} \\
                a_{21} & a_{22} & a_{23} & a_{21} & a_{22} \\
                a_{31} & a_{32} & a_{33} & a_{31} & a_{32} \\
                };

                \path ($(sarrus-1-1.north west)-(0.5em,0)$) edge ($(sarrus-3-1.south west)-(0.5em,0)$)
                    ($(sarrus-1-3.north east)+(0.5em,0)$) edge ($(sarrus-3-3.south east)+(0.5em,0)$)
                    (sarrus-1-1)                          edge            (sarrus-2-2)
                    (sarrus-2-2)                          edge[->]        (sarrus-3-3)
                    (sarrus-1-2)                          edge            (sarrus-2-3)
                    (sarrus-2-3)                          edge[->]        (sarrus-3-4)
                    (sarrus-1-3)                          edge            (sarrus-2-4)
                    (sarrus-2-4)                          edge[->]        (sarrus-3-5)
                    (sarrus-3-1)                          edge[dashed]    (sarrus-2-2)
                    (sarrus-2-2)                          edge[->,dashed] (sarrus-1-3)
                    (sarrus-3-2)                          edge[dashed]    (sarrus-2-3)
                    (sarrus-2-3)                          edge[->,dashed] (sarrus-1-4)
                    (sarrus-3-3)                          edge[dashed]    (sarrus-2-4)
                    (sarrus-2-4)                          edge[->,dashed] (sarrus-1-5);

                \foreach \c in {1,2,3} {\node[anchor=south] at (sarrus-1-\c.north) {$+$};};
                \foreach \c in {1,2,3} {\node[anchor=north] at (sarrus-3-\c.south) {$-$};};
            \end{tikzpicture}
            \]
            \item Somar os produtos das diagonais principais e somar com o 
            negativo dos produtos das diagonais secundárias:
            \[
            \left\{
            \begin{array}{l}
            a = (a_{11} \cdot a_{22} \cdot a_{33}) + (a_{12} \cdot a_{23} \cdot 
            a_{31}) + (a_{13} \cdot a_{21} \cdot a_{32}) \\
            b = - 
            ((a_{31} \cdot a_{22} \cdot a_{13}) + (a_{32} \cdot a_{23} \cdot 
            a_{11}) + (a_{33} \cdot a_{21} \cdot a_{12}))
            \end{array}
             \right.
            \]
            \item O determinante é igual a soma das diagonais principais e 
            secundárias.
        \end{enumerate}
        \item Ex.: Encontre o resultado do seguinte determinante:
        \[
        det(A_3) = 
        \begin{vmatrix}
            3 & 5 & 0 \\
            2 & 1 & 9 \\
            0 & 3 & 7 
        \end{vmatrix}
        \]
        \begin{itemize}[label=R:]
            \item 
            \[
            \begin{tikzpicture}[>=stealth]
                \matrix [%
                matrix of math nodes,
                column sep=1em,
                row sep=1em
                ] (sarrus) {%
                3 & 5 & 0 & 3 & 5\\
                2 & 1 & 9 & 2 & 1\\
                0 & 3 & 7 & 0 & 3 \\
                };

                \path ($(sarrus-1-1.north west)-(0.5em,0)$) edge ($(sarrus-3-1.south west)-(0.5em,0)$)
                    ($(sarrus-1-3.north east)+(0.5em,0)$) edge ($(sarrus-3-3.south east)+(0.5em,0)$)
                    (sarrus-1-1)                          edge            (sarrus-2-2)
                    (sarrus-2-2)                          edge[->]        (sarrus-3-3)
                    (sarrus-1-2)                          edge            (sarrus-2-3)
                    (sarrus-2-3)                          edge[->]        (sarrus-3-4)
                    (sarrus-1-3)                          edge            (sarrus-2-4)
                    (sarrus-2-4)                          edge[->]        (sarrus-3-5)
                    (sarrus-3-1)                          edge[dashed]    (sarrus-2-2)
                    (sarrus-2-2)                          edge[->,dashed] (sarrus-1-3)
                    (sarrus-3-2)                          edge[dashed]    (sarrus-2-3)
                    (sarrus-2-3)                          edge[->,dashed] (sarrus-1-4)
                    (sarrus-3-3)                          edge[dashed]    (sarrus-2-4)
                    (sarrus-2-4)                          edge[->,dashed] (sarrus-1-5);

                \foreach \c in {1,2,3} {\node[anchor=south] at (sarrus-1-\c.north) {$+$};};
                \foreach \c in {1,2,3} {\node[anchor=north] at (sarrus-3-\c.south) {$-$};};
            \end{tikzpicture}
            \]
            \[
            (3\cdot1\cdot7 + 5\cdot9\cdot0 + 0\cdot2\cdot3) 
            -
            (0\cdot1\cdot0 + 3\cdot9\cdot3 + 7\cdot2\cdot5)
            \]
            \[\equiv\]
            \[(21 + 0 + 0) - (0 + 81 + 70) \equiv 21 - 151 = -130\]

            O determinante de $A_3 = -130$ 
        \end{itemize}
    \end{itemize}
    \subsection{Teorema de Laplace}
    \begin{itemize}
        \item Além da Regra de Sarrus, como dito anteriormente, temos o Teorema 
        de Laplace, teorema este, que diferentemente da Regra de Sarrus, pode 
        ser usado para as matrizes de todas as ordens.
        \item Por mais que esse método possa ser usado para qualquer ordem, 
        ele é mais frequentemente usado em matrizes de ordem superior ou igual 
        a 4.
        \item Este teorema é bem básico, e consiste de dois passos:
        \begin{enumerate}
            \item Selecionar uma fila qualquer, com preferência para a fila 
            com a maior quantia de zeros, para facilitar os cálculos
            \item Somar os produtos dos numeros da fila selecionada pelos seus 
            respectivos cofatores
        \end{enumerate}
    \end{itemize}
    \subsubsection{Calculo dos Cofatores}
    \begin{itemize}
        \item O cofator de uma matriz de ordem superior a dois é definido como:
        \[A_{ij}=(-1)^{i+j} \cdot D_{ij}\]
        Onde:
        \begin{itemize}
            \item $A_{ij}$: Cofator de um elemento $a_{ij}$
            \item $i$: linha onde se encontra o elemento
            \item $j$: coluna onde se encontra o elemento
            \item $D_{ij}$: Determinante da matriz resultante da exclusão da 
            linha $i$ e coluna $j$
        \end{itemize}
        \item Ex.: Calculando o cofator do item $a_{12}$ da matriz $A$:
        \[A =
        \begin{vmatrix}
            3 & 5 & 0 \\
            2 & 1 & 9 \\
            0 & 3 & 7 
        \end{vmatrix}
        \]
        Com o elemento $a_{12}$, temos o cofator $A_{12} = (-1)^{1+2} \cdot 
        D_{12}$, onde $D_{12}$ é representado por:
        \[
        D_{12} = 
        \begin{vmatrix}
            2 & 9 \\
            0 & 7
        \end{vmatrix}
        = 2\cdot7 - 0\cdot 9 = 14
        \]
        Logo, o cofator do elemento $a_{12}$ é $A_{12} = (-1)^3 \cdot 14 = -14$
    \end{itemize}
    \subsubsection{Calculo do Determinante com o Teorema de Laplace}
    \begin{itemize}
        \item Calcule o determinante de $A_4$:
        \[ det(A) = 
        \begin{vmatrix}
            3 & 5 & 0 & 2 \\
            0 & 1 & 6 & 0 \\
            0 & 3 & 7 & 1 \\
            1 & 8 & 5 & 1
        \end{vmatrix}
        \]
        
        Para realizar o calculo do determinante com o Teorema de Laplace 
        precisamos selecionar alguma fila, dando preferência para a fila com 
        maior número de zeros.

        Nesse caso, podemos escolher ou a fila 2, ou a coluna 1 por possuirem 
        dois zeros, mas vamos escolher a coluna 1.
        
        A partir da coluna 1 temos: 
        \[det(A)=3 \cdot C_{11} + 0 \cdot C_{21} + 0 \cdot C_{31} + 1 \cdot C_{41}\]
        Onde $C_{11}$, $C_{21}$, $C_{31}$ e $C_{41}$ são os cofatores.

        Note que precisaremos calcular apenas dois cofatores, visto que dois 
        estão sendo multiplicados por zero
        
        \[
        \left\{
        \begin{array}{l}
            \vspace{8pt}
            C_{11} = (-1)^{i+j} \cdot D_{11} \equiv (-1)^{1+1}\cdot D_{11} \\
            \vspace{8pt}
            D_{11} = 
            \begin{vmatrix}
                1 & 6 & 0 \\
                3 & 7 & 1 \\
                8 & 5 & 1
            \end{vmatrix}
            \rightarrow
            \begin{vmatrix}
                1 & 6 & 0 \\
                3 & 7 & 1 \\
                8 & 5 & 1
            \end{vmatrix}
            \begin{matrix}
                1 & 6 \\
                3 & 7 \\
                8 & 5 
            \end{matrix} \\
            \vspace{8pt}
            D_{11} = 7 + 6\cdot 8 - 5 - 3\cdot 6 = 7 + 48 - 5 - 18 = 32 \\
            C_{11} = 1 \cdot 32 = 32
        \end{array}
        \right.
        \]

        \[
        \left\{
        \begin{array}{l}
            \vspace{8pt}
            C_{41} = (-1)^{i+j} \cdot D_{41} \equiv (-1)^{4+1}\cdot D_{41} \\
            
            \vspace{8pt}
            D_{41}
            \begin{vmatrix}
                5 & 0 & 2 \\
                1 & 6 & 0 \\
                3 & 7 & 1 
            \end{vmatrix}
            \rightarrow
            \begin{vmatrix}
                5 & 0 & 2 \\
                1 & 6 & 0 \\
                3 & 7 & 1 
            \end{vmatrix}
            \begin{matrix}
                5 & 0 \\
                1 & 6 \\
                3 & 7 
            \end{matrix} \\
            \vspace{8pt}
            D_{41} = 5\cdot 6 + 2 \cdot 7 - 2\cdot6\cdot3 = 30 + 14 - 36 = 8 \\
            C_{41} = (-1)\cdot 8 = -8
            
        \end{array}
        \right.
        \]
        Assim, o determinante de $A$ é igual a: $det(A) = 3 \cdot 32 + 0 + 0 + 
        1 \cdot (-8) = 96 - 8 = 88$
        
    \end{itemize}



\end{document}
